% ===================================================================
% 2. Related Work
% ===================================================================
\section{Related Work}
\label{sec:related}

\subsection{2D Object Detection}
\label{sec:related_detection}

Modern object detection is built on deep learning. The YOLO (You Only Look Once) family~\cite{ultralytics2024yolo11} is the go-to choice for real-time detection, treating the problem as a single regression pass. In our pipeline, we use YOLO11 in its Nano variant to keep the computational cost of the preprocessing stage low while still maintaining high recall. The detector's job here is to produce tight bounding box crops that separate the object from the background --- a necessary first step before pose estimation.

\subsection{6D Pose Estimation from RGB}
\label{sec:related_rgb_pose}

Early deep learning methods for 6D pose estimation directly predict the 3D translation and rotation from RGB images. PoseCNN~\cite{xiang2018posecnn} was one of the first end-to-end architectures to localize objects and estimate poses through a regression head. To avoid the difficulty of predicting translation directly, PVNet~\cite{peng2019pvnet} predicts 2D keypoint locations using pixel-wise voting and then recovers the 6D pose with a PnP solver, which makes it robust to occlusion. More recently, GDR-Net~\cite{wang2021gdrnet} bridges the gap between direct regression and PnP-based methods by using dense 2D--3D correspondences in a fully differentiable way. Despite their strengths, all RGB-only approaches share a fundamental problem: \emph{depth ambiguity}. Predicting the Z-axis translation depends heavily on the apparent size of the bounding box, an approximation that breaks down badly for elongated or asymmetric objects seen from different angles --- as we show in Sec.~\ref{sec:exp_results}.

\subsection{RGB-D Fusion for Pose Estimation}
\label{sec:related_rgbd}

To get around the limitations of using only color images, researchers have turned to depth sensors for direct geometric information. DenseFusion~\cite{wang2019densefusion} introduced a pixel-wise dense fusion network that extracts features from RGB images and 3D point clouds separately, then fuses them through per-pixel concatenation; each observed 3D point predicts a pose relative to itself, which naturally grounds the translation estimate in the measured geometry. Later keypoint-based RGB-D methods improved accuracy further: PVN3D~\cite{he2020pvn3d} votes for 3D keypoints from the fused point cloud, and FFB6D~\cite{he2021ffb6d} adds bidirectional fusion at every encoder layer. Both, however, supplement the scarce real LineMod training data with tens of thousands of rendered synthetic images per object, and rely on dense per-point processing.

Our work takes a different path with a ``Global Fusion'' approach: we extract a single global feature vector from each modality through independent backbones and fuse them by concatenation in a shared latent space. We keep the key geometric idea from per-point methods --- predictions should be expressed relative to observed 3D evidence --- but apply it just once, globally, through a single back-projected anchor point. This keeps the architecture lightweight and the training data real-only.

\subsection{Rotation Representations in Neural Networks}
\label{sec:related_rotation}

Representing 3D rotations in neural networks is a well-known challenge. Euler angles suffer from Gimbal Lock, and quaternions need a unit-length constraint plus they have \emph{antipodal ambiguity} ($\bq$ and $-\bq$ represent the same rotation). Naively regressing a $3\times 3$ rotation matrix (9 values) from a fully connected layer almost never produces orthogonal matrices, resulting in physically invalid transformations.

Zhou \etal~\cite{zhou2019continuity} showed that any rotation representation in $\R^n$ with $n < 5$ is necessarily discontinuous, and proposed a 6D continuous representation that converts two unconstrained 3D vectors into a valid $\SO$ matrix through a differentiable Gram-Schmidt process. This makes gradient flow during training much more stable. We use this representation in our Phase~4 model.

